\chapter{Offline-to-Online with the Latent World Model}
\label{chap:offline-online}

\section{Motivation}
\label{sec:oo-motivation}

The hierarchical pipeline developed in Chapters \ref{chap:methodology} and \ref{ch:wgsp} has taken our offline goal-conditioned agent about as far as the static dataset allows. The latent value function is calibrated; subgoals proposed by the high-level policy are well grounded by the world model; the low-level flow-matching actor handles the multi-modal action distributions of the play data without averaging across modes. What we have \emph{not} done is use the world model itself in the training loop --- it has so far only been a frozen latent space in which we plan, never a source of synthetic transitions for the critic. A natural next question is whether we can close that gap by treating the world model as an inexhaustible environment on which to continue training the offline policy.

\paragraph{From goal-conditioned to single-task.} The goal-conditioned setting is the natural home for hierarchical planning, but it is also the setting in which we have the least chance of getting a learned environment right. A goal-conditioned policy must visit a state distribution that spans the entire reachable workspace, and rolling that policy out under a learned predictor compounds error precisely along the trajectories where the model was never validated --- a goal sampled from the far side of the workspace forces the latent off the offline data's support after only a handful of WM steps. For the rest of this chapter we therefore reduce the problem to a \emph{single-task} setting: the \texttt{visual-cube-single-singletask-task1-v0} variant of OGBench, in which the goal image is fixed across episodes and the offline dataset is relabelled with the task's dense reward. Two consequences make this a useful regime to probe in isolation. First, the support the policy has to be correct on collapses to the cube-to-goal trajectories alone, narrowing both the policy's and the world model's effective operating region. Second, this is exactly the regime in which OGBench reports a privileged-state IMPALA baseline of $92.8\%$ \cite{espeholt2018impala, ogbench_park2025}, against which our offline JEPA policy plateaus at $81.6\%$; the gap is large enough that any successful method is unambiguous, and we can compare directly with the same reference numbers reported by every other paper that uses this benchmark.

\paragraph{The hypothesis.} If the LeWM predictor $f_\psi$ is accurate enough to ground subgoal planning, it ought to be accurate enough to produce additional, on-policy transitions for the agent to learn from. Wrap the predictor as a Gymnasium environment, swap it in for the real OGBench env, and continue training the offline FQL agent (\S\ref{sec:fql-bg}) using its standard losses on transitions generated by the world model. This is, in spirit, the offline-to-online paradigm of \S\ref{sec:offline-online-bg} with the learned predictor in place of additional real interaction.

\paragraph{What this chapter establishes.} Every variant of that idea we tried failed to beat its corresponding offline baseline, often by a wide margin. The bulk of the chapter is the diagnostic chain that explains why. We isolate four distinct failure modes --- compounding prediction error, $Q$-function corruption by WM-generated transitions, a termination-threshold pathology induced by WM fine-tuning, and a model-exploitation trap in joint training --- and we rule out the obvious confounder (insufficient latent expressivity) using a positive control. The synthesis in \S\ref{sec:oo-synthesis} extracts the conceptual lesson that motivates the inference-time pivot of Chapter \ref{chap:mpc-distill}: the WM cannot serve as a source of training-time supervision on this benchmark, but nothing in the failure modes rules out using it at inference for short-horizon action ranking.

\section{The latents are sufficient: a positive control}
\label{sec:oo-b2}

Before we attribute the failures of \S\S\ref{sec:oo-exp1}--\ref{sec:oo-exp8} to the world model, we have to first rule out the alternative explanation that the encoder $E_\xi$ simply does not retain enough of the scene to drive a successful policy. The cleanest test is to take everything in our pipeline except the predictor and let the agent learn against true dynamics in the real environment. We refer to this control as \emph{Latent-Obs Online RL}: the encoder is in the loop, the predictor is not.

\paragraph{Setup.} A thin observation wrapper sits between the OGBench env and the agent. Each pixel observation $o_t \in \mathbb{R}^{H\times W\times 3}$ from the real env is passed through the frozen LeJEPA encoder $E_\xi$ and only its $192$-dimensional latent $z_t = E_\xi(o_t)$ is exposed to the agent. The dynamics, rewards, and termination conditions are entirely real --- the only difference from a standard offline-to-online setup is that the observation space is the WM's encoder output rather than raw pixels. The agent itself is unchanged from the action-chunked FQL agent introduced in \S\ref{sec:fql-bg} and \S\ref{sec:hiql-flow}, with $h=5$ chunked primitive actions, the BC-flow + one-step distillation actor, and the chunked TD critic. Training runs $500\,000$ offline gradient steps on the JEPA-encoded play dataset followed by $500\,000$ steps of online interaction in which fresh real-env transitions are added to the replay buffer.

\paragraph{Result.} The agent reaches a success rate of $1.00$ at the end of the online phase, evaluated on $50$ episodes drawn from random initial states.

\paragraph{What this rules out.} The latents emitted by $E_\xi$ are sufficient, on their own, for a standard offline-to-online actor-critic to solve the task perfectly when paired with true dynamics. Any failure of the WM-environment experiments described in the rest of the chapter therefore cannot be a failure of the encoder: it must be a failure of the predictor $f_\psi$, of the supervision signal derived from it, or of the interaction between the two. We will return to this control whenever a failure mode is being diagnosed; the agent in the next several sections is exactly the agent in this section, only with $f_\psi$ providing the next state in place of the real dynamics.

\section{Common setup for the WM-environment experiments}
\label{sec:oo-setup}

All experiments in \S\S\ref{sec:oo-exp1}--\ref{sec:oo-exp8} share the same agent and the same offline dataset; they differ only in how the world model is connected to the training loop. We collect the shared elements here.

\paragraph{The WM-environment wrapper.} We wrap the frozen LeWM predictor $f_\psi$ from \S\ref{sec:lewm-bg} as a Gymnasium environment $\widetilde{\mathcal{E}} = (\mathcal{Z}, \mathcal{A}, f_\psi, \tilde r, \tilde T)$. The state space is the $192$-dimensional latent space $\mathcal{Z}$; the action space is $\mathcal{A} = [-1, 1]^{25}$, where each $25$-dimensional ``chunk'' is the row-major concatenation of $5$ consecutive $5$-D end-effector commands. Each call to $\widetilde{\mathcal{E}}.\text{step}$ thus consumes one chunk and advances the latent by one WM step, which corresponds to $5$ primitive real-env steps. Episodes terminate at most $\tilde T = 40$ WM steps after reset (equivalent to $\approx 200$ real-env steps), and reset draws the initial latent uniformly from a pool extracted from the offline dataset.

\paragraph{Reward shape.} We use one of two reward functions throughout the chapter:
\begin{align}
    \tilde r_t^{\text{sparse}} &= \mathbb{1}\!\left[\,\|z_t - z_g\|_2 < \tau_{\text{done}}\,\right], \\
    \tilde r_t^{\text{dense}}  &= -\frac{\|z_t - z_g\|_2}{s}, \qquad s = 10,
\end{align}
where $z_g$ is the latent encoding of the goal image that the real OGBench env returns at reset. Whenever the online phase uses $\tilde r^{\text{dense}}$, the offline replay buffer is relabelled with the same dense reward function before training starts, so that the critic does not have to fit two reward distributions simultaneously.

\paragraph{Why action chunks of length 5.} The choice of $5$-step chunks is not incidental: it is forced by the world model, which was trained to predict latent transitions over $5$ primitive actions at a time and has no notion of single-step dynamics. The same chunk length is what the agent is supervised on in the offline phase via the chunked Bellman target of \S\ref{sec:qchunking-bg}, so the policy and the world model agree on the temporal granularity of an action. This is the cleanest demonstration in the thesis of why action chunking is the right tool when an offline policy has to interface with a world model: the chunk length is fixed by the model, and the policy and critic must operate at the same scale to make use of it \cite{li2025reinforcement}.

\paragraph{Agent.} The agent is the action-chunked Flow Q-Learning agent described in \S\ref{sec:fql-bg}. Its critic is trained with the chunked TD target of Eq.~\eqref{eq:qchunking-target}; its actor combines the behaviour-cloning flow loss of Eq.~\eqref{eq:bc-flow-bg} with the one-step distillation objective of Eq.~\eqref{eq:fql-actor}. We use $\gamma = 0.99$, target-network Polyak rate $\tau = 0.005$, two critic heads aggregated by mean, and a BC coefficient of $\alpha = 100$ on the distillation head. Nothing in the WM-environment experiments modifies any of these losses: every section in the rest of the chapter changes only the source of the transitions $(z, \mathbf{a}, r, z')$ that the agent trains on.

\section{Plain WM-environment training}
\label{sec:oo-exp1}

The simplest possible setup: replace the real OGBench env entirely with the WM environment, use sparse rewards, and let the agent run. Each WM episode starts from a uniformly sampled latent in the initial-state pool, runs for at most $40$ WM steps with $\tau_{\text{done}}=14$, and feeds the FQL agent its standard chunked transitions. This run starts from the $61.6\%$ offline checkpoint trained with the unfinetuned LeJEPA WM.

\paragraph{Hypothesis.} The WM environment generates an essentially unbounded supply of fresh on-policy transitions in latent space. If the predictor is faithful enough, the additional coverage should let the policy continue improving in exactly the same way that online interaction with the real env does (cf.\ \S\ref{sec:oo-b2}).

\paragraph{Result.} Real-env success rate after $500\,000$ online steps: $20$--$30\%$ across three seeds.

\paragraph{Failure analysis.} The world model was trained with a one-step prediction loss, and its per-step error in cosine distance is small --- on the order of $8$--$9\%$ in the play data (verified in Chapter \ref{chap:experiments}). It is, however, unbiased: there is nothing in the loss that systematically keeps a rollout near the offline support. Over $40$ WM steps the cumulative drift is large, and the trajectories the agent ends up training on visit latents that the predictor itself was never trained to model --- the latent-space analogue of a robot wandering into a room that does not exist. The critic, fit to TD targets computed in those drifted regions, learns to extrapolate Q-values into hallucinated states; the actor then exploits the extrapolation by selecting actions that maximise the corrupted critic. On the real env, that policy executes actions that no longer correspond to any meaningful behaviour, and performance collapses well below the offline baseline. The contrast with \S\ref{sec:oo-b2} is sharp: the same agent, with the same encoder and the same offline initialisation, reaches $100\%$ when paired with true dynamics and $20$--$30\%$ when paired with $f_\psi$. The locus of the failure is the predictor.

\section{Anchored restart rollouts}
\label{sec:oo-exp2}

Compounding prediction error is the most obvious symptom of \S\ref{sec:oo-exp1}, and the most direct fix is to cap the length of any unconstrained rollout. We do this with an \emph{anchored restart rollout} (sometimes called Variant B in the diagnostic chain): every $L_b$ WM steps the current latent is matched to its nearest neighbour in the offline pool $\mathcal{Z}_{\text{off}}$, and if the match is close enough the rollout is truncated and restarted from the anchor on the next reset.

\paragraph{Anchor mechanism.} The nearest-neighbour query against $\mathcal{Z}_{\text{off}}$ is implemented via a GPU-resident FAISS index (\S\ref{sec:faiss-bg}), which makes the query effectively free at WM-rollout speeds. Define the threshold-projected next-state operator
\begin{equation}
    \Pi^{\epsilon}_{\mathcal{Z}_{\text{off}}}(z) \;=\;
    \begin{cases}
        \hat z \;=\; \arg\min_{z' \in \mathcal{Z}_{\text{off}}} \|z - z'\|_2 &
        \text{if } \cos(z, \hat z) \ge \tau_{\cos}\;\text{and}\;\|z - \hat z\|_2 \le \tau_{L_2},\\[4pt]
        z & \text{otherwise,}
    \end{cases}
    \label{eq:oo-anchor-proj}
\end{equation}
applied at end-of-branch boundaries every $L_b = 3$ WM steps, with $\tau_{\cos} = 0.95$ and $\tau_{L_2} = 4.0$. If the nearest neighbour passes the threshold, the rollout is truncated and the next \texttt{reset()} starts from $\hat z$; if not, the rollout is abandoned and a fresh latent is drawn from the initial-state pool. Importantly, the transition tuple that is appended to the replay buffer always uses the raw WM output as $z_{t+1}$ --- there is no fake snap-back edge in the learned trajectory.

\paragraph{Three variants.} We ran the anchored protocol in three configurations of growing ambition:
\begin{itemize}
    \item \emph{Sparse-reward, base WM} (the original LeJEPA predictor, $\tau_{\text{done}} = 10$, starting from the $61.6\%$ offline checkpoint).
    \item \emph{Dense-reward, partially fine-tuned WM} (the base WM with its predictor head fine-tuned on the offline play dataset, $\tau_{\text{done}} = 15$).
    \item \emph{Dense-reward, fully fine-tuned WM} (both predictor and encoder fine-tuned on the offline play data, $\tau_{\text{done}} = 15$, starting from the $81.6\%$ checkpoint).
\end{itemize}
Real-env success rates respectively: $13$--$22\%$, $8$--$54\%$ (very high variance), and $15.2\%$.

\paragraph{Failure analysis, part one: anchoring fixes drift but not Q-corruption.} Anchoring does keep the per-branch rollout bounded, so the most extreme drift-induced hallucinations of \S\ref{sec:oo-exp1} are avoided. The critic, however, is still trained on sparse-reward signals computed against the WM's notion of ``goal reached''. That notion does not reliably agree with the real env's: a rollout that the WM thinks has reached the goal --- and therefore receives a reward of $1$ --- may correspond to a real-env trajectory that has not. Q-values learnt against an unreliable success indicator are themselves unreliable, even when the underlying latent trajectory is well anchored.

\paragraph{Failure analysis, part two: a termination-threshold pathology.} The full fine-tune variant fails for a more specific reason that is worth dwelling on, because it illustrates the brittleness of any reward function derived from the WM's geometry. The base LeJEPA predictor produces a typical agent-to-goal distance of $\approx 25$ in the offline data; the partially fine-tuned predictor produces $\approx 19$; the fully fine-tuned predictor produces $\approx 11$. With $\tau_{\text{done}} = 15$ --- a threshold calibrated for the base WM --- every starting latent in the fully fine-tuned WM is already inside the success region, so every WM episode terminates at step $1$ with the success flag set. The agent is supervised throughout the entire $500\,000$ online steps on degenerate one-step episodes that carry no horizon, no credit assignment, and no exploration; the offline policy is actively corrupted by training on this signal, and the $81.6\%$ baseline drops to $15.2\%$. The local fix is to recalibrate $\tau_{\text{done}}$ per WM, but the underlying lesson is more general: any reward shaping built on top of the WM's latent geometry is sensitive to its calibration, and our fine-tuning pipeline produces predictors whose geometries differ enough to make uniform thresholds unsafe.

\section{Joint training of the world model and policy}
\label{sec:oo-exp3}

The failures of \S\S\ref{sec:oo-exp1}--\ref{sec:oo-exp2} suggest that a frozen WM may simply be the wrong asset to train on, and that updating the predictor against the same value signal the policy uses might align it with the task. We test that idea in this section, but before doing so we have to be precise about \emph{which} WM is being modified: the rest of the chapter refers to several distinct fine-tuning stages, and the distinction matters.

\paragraph{The four world models in this chapter.} We use four predictors, all built from the same LeJEPA architecture but trained on different data and with different gradient pathways:
\begin{itemize}
    \item \emph{Base WM}: trained on the OGBench \texttt{visual-cube-single-play-v0} dataset, which consists of undirected play-style exploration on the single-cube setup. This is the predictor used by all goal-conditioned methods in earlier chapters; it has never seen the task-relabelled offline dataset.
    \item \emph{Predictor-finetuned WM}: the base WM with only its predictor head fine-tuned on the offline play dataset. The encoder is frozen at its base weights.
    \item \emph{Fully fine-tuned WM}: the base WM with both encoder and predictor fine-tuned on the offline play dataset. This is the WM that produces the latent-distance compression discussed in the previous section.
    \item \emph{Joint-trained WM}: the WM in this section, trained alongside the policy and critic during the offline phase with the value-aware loss described below.
\end{itemize}
Whenever ambiguous, we name the WM explicitly in the experiment description.

\paragraph{Joint-training loss.} The joint-trained WM has its own AdamW optimiser and receives gradient from a two-term loss applied during the offline phase:
\begin{align}
    \mathcal{L}_{\text{WM}}(\psi) &= \alpha\,\mathcal{L}_{\text{pred}}(\psi) + \beta(t)\,\mathcal{L}_{\text{value}}(\psi),
    \label{eq:oo-joint}\\
    \mathcal{L}_{\text{pred}}(\psi) &= \mathbb{E}\!\left[\,
        \frac{\|f_\psi(z, a) - z'\|_2^2}{\|z'\|_2^2 + \varepsilon}\,\right],
    \label{eq:oo-pred}\\
    \mathcal{L}_{\text{value}}(\psi) &= \mathbb{E}\!\left[\Big(
        Q_{\bar\theta}(z, a) - \big(r + \gamma\, m\, Q_{\bar\theta}(f_\psi(z, a),\, a')\big)\Big)^2\right],
    \quad a' \sim \pi_\theta\big(\,\cdot\mid \mathrm{sg}[f_\psi(z, a)]\big).
    \label{eq:oo-value}
\end{align}
The intent of $\mathcal{L}_{\text{value}}$ is to encourage the WM to produce predictions that are \emph{self-consistent} with the critic --- i.e., predicted next states whose target-critic value satisfies the Bellman equation against the real reward. The stop-gradient $\mathrm{sg}[\cdot]$ keeps the policy from receiving updates through the WM; the target critic $Q_{\bar\theta}$ keeps the critic itself from becoming a moving target during early training. The relative weight $\beta(t)$ ramps linearly from $0$ to $\beta_{\max} = 0.1$ over the first $\beta_{\text{ramp}} = 50\,000$ steps, giving the critic time to settle before the WM starts to listen to it; $\alpha = 1$. The actor and critic losses are unchanged from Eqs.~\eqref{eq:bc-flow-bg}--\eqref{eq:fql-critic}.

\paragraph{Result.} After offline joint training, the joint policy reaches $79.6\%$ on its own; running an anchored, dense-reward online phase using the joint WM as the environment and the joint policy as the starting point yields $9.6\%$.

\paragraph{Diagnostic decomposition.} To localise the failure, we ran two follow-up experiments that pair each joint-trained component with its non-joint counterpart in the online phase:
\begin{center}
\begin{tabular}{lll r}
\toprule
Run & Policy used online & WM used online & SR (\%) \\
\midrule
\emph{base policy + joint WM}  & base ($81.6\%$) & joint WM &  $7.2$ \\
\emph{joint policy + base WM}  & joint ($79.6\%$) & base WM & $30.0$ \\
\bottomrule
\end{tabular}
\end{center}
The joint policy is no worse than other failed runs of this chapter when paired with the base WM ($30\%$, in line with the original $20$--$30\%$ of \S\ref{sec:oo-exp1}). The joint WM destroys performance regardless of which policy it is paired with ($7.2\%$ with the offline-only $81.6\%$ baseline). The locus of the failure is the predictor, not the policy.

\paragraph{Failure analysis: model exploitation.} The pattern is the canonical model-exploitation failure of model-based RL. The value-residual loss $\mathcal{L}_{\text{value}}$ measures how well the target-critic agrees on both sides of a one-step Bellman equation, but it does not anchor $f_\psi$ to any particular ground-truth dynamics: nothing in Eq.~\eqref{eq:oo-value} prevents the predictor from outputting a single high-value latent $z_{\text{good}}$ irrespective of its inputs, so long as $Q_{\bar\theta}(z_{\text{good}}, \cdot)$ is large. The prediction loss $\mathcal{L}_{\text{pred}}$ is the only counterweight, and even at its final relative weight of $10\times \mathcal{L}_{\text{value}}$ it is pulling in a much harder direction --- predicting the correct next latent for arbitrary inputs --- than the value loss --- producing one high-value latent for everything. The gradient field is structurally biased toward the easier minimum. MOPO \cite{yu2020mopomodelbasedofflinepolicy} and COMBO \cite{yu2021combo} address exactly this concern in the supervised-WM-then-pessimism style: keep the model frozen and penalise the TD target where the model is uncertain. \S\ref{sec:oo-exp45} tests that family of methods directly.

\section{Uncertainty-penalised rewards}
\label{sec:oo-exp45}

A standard remedy for model exploitation is to keep the WM frozen but penalise predicted transitions by some estimate of model uncertainty, so the agent learns to avoid regions where the predictor is least reliable \cite{yu2020mopomodelbasedofflinepolicy, yu2021combo}. We tested two variants of this idea inside the anchored, dense-reward WM-environment pipeline, both starting from the $81.6\%$ offline baseline. The penalised reward is
\begin{equation}
    \tilde r'_t \;=\; \tilde r_t \;-\; \lambda\, u(z_{t+1}, a_t),
    \label{eq:oo-mopo-reward}
\end{equation}
with two choices of uncertainty estimator $u$.

\paragraph{Ensemble disagreement.} Three independently seeded copies of the fully fine-tuned WM, $\{f_\psi^{(k)}\}_{k=1}^3$, are run in parallel at every step; the per-step uncertainty is the $\ell_2$ norm of the per-element standard deviation of their predictions:
\begin{equation}
    u_{\text{ens}}(z, a) \;=\; \Big\|\,\sigma_k\!\big[f_\psi^{(k)}(z, a)\big]\,\Big\|_2,
    \label{eq:oo-mopo-ens}
\end{equation}
with $\lambda = 0.6$. Result: $21.2\%$. This is a uniform downward offset on rewards across the WM's operating regime --- the three ensemble members agree on most rollouts because they were all fine-tuned on the same offline data --- so the penalty is less destructive than the alternatives below but also does not discriminate between good and bad actions in a useful way.

\paragraph{Nearest-neighbour cosine distance.} A single WM, with uncertainty estimated as the cosine distance from each predicted latent to its nearest neighbour in the offline pool:
\begin{equation}
    u_{\text{nn}}(z') \;=\; 1 \;-\; \max_{z\in\mathcal{Z}_{\text{off}}}\;
        \frac{\langle z, z'\rangle}{\|z\|\,\|z'\|},
    \label{eq:oo-mopo-nn}
\end{equation}
with $\lambda = 0.07$. The NN query reuses the FAISS index from \S\ref{sec:oo-exp2}. Result: $2.4\%$.

\paragraph{Failure analysis.} The NN-distance result is striking because it is worse than the unpenalised baseline, and the reason is structural rather than a tuning issue. The play data was collected by undirected exploration of the workspace, so the latents it produces are concentrated in the workspace's interior rather than along the goal-directed trajectories that solve the task. An NN-distance penalty therefore measures \emph{how far this state is from the play prior}, not how unreliable the WM's prediction at this state actually is. Successful task-solving actions push the cube toward the goal latent --- which is exactly the part of the latent space the play data does not cover --- so the penalty is largest precisely on the actions we want the policy to take. The Q-function learns to rank good actions below noisy in-distribution actions, and the policy collapses. The ensemble variant escapes this specific trap because the three predictors all share the play-data bias, so their disagreement does not point in the same direction as the bias; but neither penalty discriminates between WM error and policy correctness in a way that actually helps, and both runs additionally inherit the degenerate one-step-episode pathology of \S\ref{sec:oo-exp2}.

\section{Differentiable WM rollouts}
\label{sec:oo-exp8}

A final option is to abandon the WM-as-environment framing entirely and use the predictor instead as a differentiable simulator, backpropagating an analytic policy gradient through a short imagined rollout into the actor parameters. This is the classical formulation of stochastic value gradients \cite{heess2015svg} and the same idea Dreamer \cite{hafner2020dreamer} uses to update its actor through imagination.

\paragraph{Rollout loss.} The actor receives an additional term
\begin{equation}
    \mathcal{L}_{\text{roll}}(\theta) \;=\;
        -\mathbb{E}_{z\sim\mathcal{D}}\!\left[\,\sum_{k=0}^{H-1}\gamma^{k}\,
            \tilde r^{\text{dense}}\!\big(z_{t+k+1}\big)\right],
    \quad
    z_{t+k+1} = f_\psi\!\big(z_{t+k},\,\pi_\theta(z_{t+k})\big),
    \label{eq:oo-rollout-loss}
\end{equation}
with $H = 3$. Both $f_\psi$ and the flow-matching actor $\pi_\theta$ are differentiable, so the gradient of the dense reward at the rollout endpoint flows backward through three WM steps and through the actor's $S = 10$ Euler integration steps into the actor parameters. The total actor loss becomes $\mathcal{L}_{\text{actor}} + \lambda_{\text{roll}}\,\mathcal{L}_{\text{roll}}$, with $\lambda_{\text{roll}}$ controlling the weight of the analytic-gradient term. In contrast with \S\ref{sec:oo-exp3}, the WM here is \emph{frozen}: the actor's gradients flow through it without modifying it.

\paragraph{Result.} For both $\lambda_{\text{roll}} \in \{0.1, 1.0\}$, real-env success rate after $500\,000$ online steps: $0.0\%$.

\paragraph{Failure analysis.} The gradient passes through a chain of $S + H = 13$ noisy compositions before reaching the dense reward, and each composition adds variance to the estimate. The flow integrator alone is a sequence of $10$ velocity-field evaluations on noisy intermediate states; cascading $3$ further WM steps on top --- each contributing $\approx 8\%$ cosine error --- amplifies the variance enough that the analytic gradient no longer reliably points in the direction of higher real-env return. At $\lambda_{\text{roll}} = 1$ the rollout term dominates the BC-flow loss and the actor abandons the behavioural distribution to chase the WM's dense reward; at $\lambda_{\text{roll}} = 0.1$ the rollout gradient is too small to influence the actor against the BC term. There is no setting of $\lambda_{\text{roll}}$ that is both informative and not destructive.

\section{Synthesis}
\label{sec:oo-synthesis}

Table~\ref{tab:oo-summary} collects every WM-based offline-to-online attempt of this chapter. None beats its corresponding offline baseline.

\begin{table}[h]
\centering
\caption{Summary of all WM-based offline-to-online attempts on
   \texttt{visual-cube-single-singletask-task1-v0}. SR is the
   real-env success rate over $250$ episodes. The latent-obs control
   row (\S\ref{sec:oo-b2}) is included for reference: with the real env
   in place of the WM, the same agent reaches $100\%$.}
\label{tab:oo-summary}
\begin{tabular}{l l c c}
\toprule
Experiment & Mechanism & Offline SR (\%) & Online SR (\%) \\
\midrule
Latent-Obs Online RL (\S\ref{sec:oo-b2})       & encoder-only control                       & 81.6 & \textbf{100.0} \\
\midrule
Plain WM-env, sparse (\S\ref{sec:oo-exp1})       & no anchoring                       & 61.6 & 20--30 \\
Anchored, sparse (\S\ref{sec:oo-exp2})           & $L_b{=}3$, $\tau_{\cos}{=}0.95$    & 61.6 & 13--22 \\
Anchored, dense, partial FT (\S\ref{sec:oo-exp2})& dense reward                       & 61.6 & 8--54  \\
Anchored, dense, full FT (\S\ref{sec:oo-exp2})   & $1$-step degenerate episodes       & 81.6 & 15.2   \\
Joint training, anchored (\S\ref{sec:oo-exp3})   & $\mathcal{L}_{\text{pred}}+\beta\mathcal{L}_{\text{value}}$ & 79.6 & 9.6    \\
\quad diag: base policy + joint WM               & isolate WM                         & 81.6 & 7.2    \\
\quad diag: joint policy + base WM               & isolate policy                     & 79.6 & 30.0   \\
Ensemble penalty (\S\ref{sec:oo-exp45})          & $\lambda{=}0.6$, $3$ WMs            & 81.6 & 21.2   \\
NN-distance penalty (\S\ref{sec:oo-exp45})       & $\lambda{=}0.07$                    & 81.6 & 2.4    \\
Differentiable rollouts (\S\ref{sec:oo-exp8})    & frozen WM, actor BPTT               & 61.6 & 0.0    \\
\bottomrule
\end{tabular}
\end{table}

\paragraph{Four failure mechanisms.} The encoder-only control of \S\ref{sec:oo-b2} rules out the encoder, so every failure above is attributable to the world-model predictor. The diagnostic chain identifies four distinct routes by which it goes wrong:
\begin{enumerate}
    \item \emph{Compounding latent-prediction error.}
      Per-step error of $\sim 8\%$ in cosine distance compounds over $40$-step rollouts. The drifted latents take both policy and critic outside the offline support
      (\S\ref{sec:oo-exp1}, \S\ref{sec:oo-exp2}).
    \item \emph{$Q$-function corruption.}
      WM-only transitions feed Bellman targets that the critic fits;
      since the WM transitions are systematically wrong, the resulting
      $Q$ ranks real-env actions incorrectly. Every online run suffers
      this regardless of anchoring or pessimism.
    \item \emph{Termination-threshold pathology.}
      Fine-tuned WMs compress latent distances. With a fixed
      $\tau_{\text{done}}$, episodes terminate at step~$1$, producing
      useless training signal (\S\ref{sec:oo-exp2}, \S\ref{sec:oo-exp45}).
    \item \emph{Model-exploitation trap.}
      When the WM is updated against a critic-derived loss it learns to
      hallucinate states that satisfy the Bellman equation rather than
      states that correspond to real dynamics (\S\ref{sec:oo-exp3}).
\end{enumerate}

\paragraph{What the failures do not rule out.} The four mechanisms above are all problems of the WM as a \emph{source of supervision} for the policy or critic: TD targets, dense reward, or analytic gradients. What none of them challenges is the weaker claim that the WM, applied for one or two steps starting from a real on-policy latent, can produce an informative \emph{ranking} of candidate action chunks under the learned critic. We have not used the WM that way in any experiment in this chapter --- the closest is the score function that gates the anchoring projection in \S\ref{sec:oo-exp2}, which is a coarse threshold rather than a learned ordering. Chapter \ref{chap:mpc-distill} takes exactly that approach: training stays fully offline on real data so the critic remains honest, and the WM is used only at inference time to expand the action search beyond what the actor would propose on its own.
